\section*{Chapter \ref{ch:nw:fx-connectivity} --- FX Connectivity}

\begin{solution}{exo:nw:fx-connectivity:1}
New York FX spot (and metals) in Equinix Secaucus NY5, with a backup in Slough LD4.2; London FX spot and NDFs in LD4.2, with a backup
in NY5. Each product trades in a single location.
\end{solution}

\begin{solution}{exo:nw:fx-connectivity:2}
Only the price levels where the firm has bilateral credit with the counterparties; the quantity at a level may include orders it
cannot trade with, so it may exceed what the firm can fill.
\end{solution}

\begin{solution}{exo:nw:fx-connectivity:3}
\qty{46.85}{\milli\second} one way; at 1.5, \qty{70.3}{\milli\second} one way and \qty{140.5}{\milli\second} round trip.
\end{solution}

\begin{solution}{exo:nw:fx-connectivity:4}
The provider is at the origin: its quote leaves for each site as soon as it has priced the move, over a path as fast as the taker's, so
the only difference is the \qty{0.1}{\milli\second} of pricing.
\end{solution}

\begin{solution}{exo:nw:fx-connectivity:5}
\qty{79.5}{\milli\second}: the network at 1.5 plus \qty{40}{\micro\second} of stages. A server in Secaucus next to the engine would
save almost all of it, but would learn of Tokyo news \qty{79.5}{\milli\second} late, so it helps only for New York-driven trading.
\end{solution}

\begin{solution}{exo:nw:fx-connectivity:6}
After a move, the quotes that have not yet moved are the best on the side the market moved away from, and they come from the providers
furthest from the move; the aggregator selects them because they look best, exactly while they are stale.
\end{solution}

\begin{solution}{exo:nw:fx-connectivity:7}
\qty{47}{\milli\second}: 99.4\% are caught; at \qty{46}{\milli\second}, 96.2\%.
\end{solution}

\begin{solution}{exo:nw:fx-connectivity:8}
Staleness is not what the check must cover: with equal paths the price check at arrival already knows the move. If the fills are toxic,
the provider's information path is slower than the takers'; a window longer than the checks need is contrary to Principle 17. The fix is
a faster information path or pricing in Tokyo.
\end{solution}

\begin{solution}{pb:nw:fx-connectivity:1}
\textbf{Part I.}
\begin{enumerate}
\item \qty{27.07}{\milli\second} and \qty{46.85}{\milli\second}.
\item \qty{40.6}{\milli\second} and \qty{70.3}{\milli\second}.
\item \qty{140.6}{\milli\second} at TY3, \qty{31.5}{\milli\second} at NY5 and NY6, \qty{111.1}{\milli\second} at SG1, and
  \qty{0.1}{\milli\second} at LD4, where the taker learns of it no sooner than the provider.
\item \qty{81.3}{\milli\second} at NY5 and NY6, \qty{31.5}{\milli\second} at TY3, \qty{8.2}{\milli\second} at SG1 and
  \qty{0.1}{\milli\second} at LD4.
\end{enumerate}
\textbf{Part II.}
\begin{enumerate}[resume]
\item None beyond the pricing time: \qty{0.1}{\milli\second}.
\item 70\% at once and 99.4\% after \qty{2}{\milli\second}.
\item \qty{185.6}{\milli\second} and \qty{45.1}{\milli\second}.
\item \qty{47}{\milli\second}.
\end{enumerate}
\textbf{Part III.}
\begin{enumerate}[resume]
\item A validity check, including sufficient available credit, and a price check against the current price available to the client.
\item No: a delay beyond what the checks need is contrary to Principle 17.
\item A faster information path from Tokyo, or pricing Tokyo moves from a server in Tokyo; disclosed, prompt checks.
\item In TY3, where the streams to the Tokyo site start: the Tokyo site's quotes would be fresh after Tokyo moves, and stale after London
  and New York ones instead.
\end{enumerate}
\textbf{Part IV.}
\begin{enumerate}[resume]
\item \emph{Named result}: a London provider's USDJPY quote is \qty{140.6}{\milli\second} stale at the Tokyo site after a Tokyo move with
  every path at 1.5 times the floor, and \qty{185.6}{\milli\second} if it learns of Tokyo moves over a backbone at 2.46; the hold that would
  neutralise it is only the pricing time with equal paths, and \qty{45.1}{\milli\second} (47 to catch 99\% with jitter) with the backbone, a
  delay the Code does not allow, so the remedy is the path, not the window.
\item Published: the sites, the Code's text, chapter~19's \omterm{def:nw:the-north-american-map:factor}{route factors}. Assumed: the factors applied to these routes, the pricing time, the
  jitter, and where the news breaks.
\item The Tokyo provider's quote fresh, the London provider's stale for \qty{140.5}{\milli\second}, the New Jersey provider's for
  \qty{158.9}{\milli\second}; it selects the stale ones on the side the market left.
\item Timestamp each quote at creation and at each venue's acknowledgement or market data echo, and compare fills with its own mid a short
  time later, per site and per origin of the move.
\item Its orders are forwarded to New York or London, so it trades on the engine's time, a crossing away, whatever its gateway.
\item They shorten the takers' paths: the provider's information path is then slower than theirs, and the hold needed grows by the
  difference, again a matter of paths, not windows.
\item In the FX rows: the venues and sites per pair, the provider's pricing sites and paths, credit, and the staleness monitoring.
\item Whether the provider learns of a move as fast as the taker's request can reach it.
\end{enumerate}
\end{solution}

\begin{solution}{iq:nw:fx-connectivity:1}
In LD4 in Slough and on the Secaucus campus, where the main engines and streams are, with a presence in TY3 for Tokyo clients and
streaming; SG1 for NDFs. Each pair matches in one place, so the servers go where the pairs you trade match and where their news breaks.
\iqlookfor{Engines per pair; the Slough and Secaucus campuses; Tokyo as access and streaming.}
\end{solution}

\begin{solution}{iq:nw:fx-connectivity:2}
A risk control for validity, including credit, and price, applied promptly; it is not a free option, a source of information for trading
against the client, or a way to add delay beyond the checks.
\iqlookfor{Principle 17; validity and price; no added delay; disclosure.}
\end{solution}

\begin{solution}{iq:nw:fx-connectivity:3}
Timestamp quotes at creation on a synchronised clock, read the venues' market data or acknowledgements with their timestamps at each site,
and compare with the moment the move that should have changed the quote reached each site by the fastest path.
\iqlookfor{Clock synchronisation; per-site measurement; the fastest path as reference.}
\end{solution}

\begin{solution}{iq:nw:fx-connectivity:4}
Whether the provider's information path changed (a route, a vendor, a switch), whether takers gained a faster path, the rejects by reason,
and the fills against the mid a few milliseconds later, per site.
\iqlookfor{Paths on both sides; markouts; reject reasons.}
\end{solution}

\begin{solution}{iq:nw:fx-connectivity:5}
It is the maximum over stale and fresh quotes, and after a move the stale ones are the best; with bilateral credit and last look, the best
shown price may also be one the client cannot trade or that will be rejected.
\iqlookfor{Selection of stale quotes; credit screening; rejects.}
\end{solution}

\begin{solution}{iq:nw:fx-connectivity:6}
Pricing engines in London, Secaucus and Tokyo, each fed by the fastest paths from the others and from the pairs' main markets; streams from
the nearest engine to each site; one credit service consistent across venues; prompt, disclosed price and validity checks; staleness and
markouts measured per site.
\iqlookfor{Pricing near the news; path parity; credit; monitoring.}
\end{solution}
